\documentclass[11pt]{article}
\usepackage[margin=0.68in]{geometry}
\usepackage{amsmath}
\usepackage{booktabs}
\usepackage[table]{xcolor}
\usepackage{graphicx}
\usepackage{float}
\usepackage{caption}
\usepackage{longtable}
\usepackage{pdflscape}
\usepackage{newtxtext,newtxmath}
\usepackage{microtype}
\usepackage[hidelinks]{hyperref}
\setlength{\parindent}{0pt}
\setlength{\parskip}{0.55em}
\pagestyle{plain}

\title{Predictable Financial Crises:\\A Public-Data Replication and Post-Publication Update}
\author{Brian Nguyen and Clara Duan\\FINM 32900}
\date{August 2026}

\begin{document}
\maketitle

\begin{abstract}
This project reconstructs the country-year data and principal results in
Greenwood, Hanson, Shleifer, and S{\o}rensen's \textit{Predictable Financial
Crises} using public data. The paper's central object is the R-Zone: a year in
which three-year debt growth is in the historical top quintile and real
asset-price growth is in the historical top tercile. We reproduce the assigned
historical tables and event-history figures, validate them against 586
published benchmarks using documented tolerances, extend the predictors and
R-Zone classifications through 2025, and investigate bank fragility, missed
crises, dynamics, and the United States around the 2023 banking stress. The
replication succeeds most clearly in recovering the sample frequencies,
cross-cell crisis gradient, and qualitative regression pattern. Exact equality
is limited chiefly by unavailable proprietary equity-price data, revised
public-source vintages, and differences in finite-sample inference.
\end{abstract}

\tableofcontents
\clearpage

\section{Nature and scope of the replication}

The original paper asks whether financial crises are predictable from the
joint behavior of credit and asset prices. We reconstruct annual observations
for the paper's 42 countries, as close as public coverage permits, beginning in
1950. Business-sector R-Zones pair business debt-to-GDP growth with real equity
price growth; household-sector R-Zones pair household debt-to-GDP growth with
real house-price growth. The historical forecast-origin sample ends in 2012 so
that crisis outcomes are observed over a common one- through four-year future
window ending in 2016.

The project has three layers. First, it recreates the historical variables,
descriptive exhibits, and country-fixed-effects probability models. Second, it
holds the historical thresholds fixed and extends the underlying predictors
and main exhibits through the latest usable observations. Third, it studies
whether aggregate bank-balance-sheet fragility and autoregressive memory add
information beyond the original borrower-side R-Zone indicators.

\section{Data sources and construction}

The crisis chronology, bank-equity crashes, bank failures, panics, and the
historical GDP series come from the Baron--Verner--Xiong replication data.
Credit is assembled from the IMF Global Debt Database, BIS Total Credit
Statistics, and the Jord\`a--Schularick--Taylor Macrohistory database. Real
equity prices use the IMF's former share-price series, JST, and OECD share
prices. Real house prices use BIS Residential Property Prices, OECD Analytical
House Prices, and JST. CPI and nominal GDP needed for real transformations come
from World Bank World Development Indicators with JST supplements.

Each three-year change is calculated within a source before sources are
spliced. This avoids constructing a growth rate across incompatible levels.
Debt growth is measured as a three-year percentage-point change in debt to GDP;
real asset-price growth is 100 times a three-year log difference. Missing
values remain missing, and a future crisis outcome is defined only when the
entire relevant future window is covered. The historical pooled Q80 debt and
T66.7 price cutoffs are frozen for every later classification.

For the post-2016 crisis update, the outcome series extends BVX's own
definition forward, so ``crisis'' means one thing in every year of the panel:
a bank equity index decline of 30\% or more together with narrative evidence
of widespread bank failures or a banking panic. No published chronology
applies these criteria past 2016, so the extension is constructed here from
published datasets: the BVX replication data and the CRSP value-weighted bank
portfolio from the Ken French Data Library, which tracks the paper's own
United States bank equity series closely over their post-war overlap. A
narrative screen of the sample for 2017--2025 yields three candidate
episodes, and the pulled data decide each label mechanically
(Table~\ref{tab:crisis-screen}). The 2023 United States episode is genuinely
borderline: the broad value-weighted index stops short of the 30\% bar
because the crash concentrated in regional banks while the largest
institutions rallied, even though large-cap bank benchmarks crossed the bar
and the failure and run evidence is unprecedented by assets. BVX's own
chronology contains a precedent in the other direction---the United States
in 1984 is counted as a crisis on narrative evidence with a smaller
aggregate equity decline---so we let the mechanical rule decide and report
the sensitivity openly.

\begin{table}[H]
\centering
\caption{\textbf{Post-2016 candidate crisis episodes under the paper's own
criteria.} Each row is a sample episode from the 2017--2025 narrative screen,
with the verdict of BVX's criteria as applied by our pipeline. The takeaway
is that no post-2016 episode qualifies as a systemic crisis onset, but the
United States in 2023 is borderline: the broad value-weighted bank index
falls short of the 30\% threshold that large-cap benchmarks crossed, while
the failure and run evidence rivals any postwar episode. An index choice,
not the data alone, decides the 2023 call.}
\label{tab:crisis-screen}
\small
\input{../_output/post_publication_crisis_screen.tex}
\end{table}

\begin{table}[H]
\centering
\caption{\textbf{Original summary statistics for the underlying predictor
data.} The table compares complete sector-specific predictor pairs in the
historical estimation years with later observations. The main lesson is that
country coverage becomes much broader after 2012, while average debt growth
and the fraction of joint upper-tail R-Zones fall. Updated estimates therefore
reflect both more countries and a different predictor distribution.}
\label{tab:overview}
\small
\input{../_output/data_overview_summary.tex}

\medskip
\begin{minipage}{0.97\linewidth}
\footnotesize
\textit{Notes:} Coverage is relative to a balanced panel of 42 countries. Debt
and price entries report mean (standard deviation). Corr. is the Pearson
correlation between the two predictors. R-Zone percentages use the frozen
historical thresholds.
\end{minipage}
\end{table}

\begin{figure}[H]
\centering
\includegraphics[width=0.91\textwidth]{../_output/data_overview_figure.jpg}
\caption{\textbf{Coverage and joint predictor distributions.} Panel A shows
the changing number of countries with a complete sector pair. Panels B and C
show later predictor observations and the frozen threshold lines. The figure
motivates two cautions for interpretation: annual global statistics do not use
a constant country count, and an R-Zone is a genuinely rare joint-tail event
rather than a year in which only debt or only prices are elevated. Scatter axes
are limited to the first and 99th percentiles for readability.}
\label{fig:overview}
\end{figure}

\clearpage
\section{Historical replication}

\subsection{Summary statistics and event histories}

The reconstructed sample closely matches the paper's core moments. Business
R-Zones account for 6.46\% of complete historical business observations,
compared with 6.0\% in the paper; household R-Zones account for 9.81\%, compared
with 10.3\%. The three-year crisis frequency conditional on an R-Zone is
36.6\% for business and 33.6\% for households in the public reconstruction,
versus 45.3\% and 36.8\% in the paper. The remaining business gap is consistent
with the equity-series substitution discussed below.

\begin{table}[H]
\centering
\caption{\textbf{Replicated Table 1: historical summary statistics.} Every
entry is recalculated from the reconstructed public-data panel using the
paper's sample logic. The close crisis means and percentile thresholds support
the overall reconstruction; the largest remaining difference is equity-price
dispersion because the proprietary GFD and Bloomberg series are unavailable.}
\label{tab:table1-historical}
\small
\input{../_output/table1_summary_stats.tex}
\end{table}

\begin{figure}[p]
\centering
\includegraphics[width=0.96\textwidth]{../_output/historical_business_rzone_timeline.png}
\caption{\textbf{Replicated Figure 1: historical event history.} Coral circles
mark BVX crisis onsets and teal crosses mark business R-Zones. Horizontal
country lines show the span with both predictors. Together with Panel B on the
following page, the figure shows that R-Zones often cluster before crises but
are neither necessary nor sufficient.}
\label{fig:figure1-historical}
\end{figure}

\begin{figure}[p]
\ContinuedFloat
\centering
\includegraphics[width=0.96\textwidth]{../_output/historical_household_rzone_timeline.png}
\caption[]{\textbf{Replicated Figure 1 continued: household event history.}
The household panel reinforces the same takeaway: some joint booms do not end
in crisis, and many crises arrive without a preceding R-Zone.}
\end{figure}

\clearpage
\subsection{Crisis probabilities across predictor cells}

The descriptive cell estimates preserve the paper's main gradient: crisis
frequencies generally rise when both price growth and debt growth are high.
Sparse cells and substituted equity data create numerical differences, but the
upper-right R-Zone remains economically distinct from the median cell.

\captionof{table}{\textbf{Replicated Table 3: crisis probabilities by debt and
price-growth quantiles.} Panels A and C show where the observations lie;
Panels B and D show one- through four-year crisis frequencies and differences
from the median cell. Gray upper-right cells identify the R-Zone. The main
lesson is a joint-tail pattern rather than a purely monotonic effect of either
predictor by itself. Significance uses conventional Driscoll--Kraay p-values.}
\label{tab:table3-historical}
\input{../_output/table3_replication.tex}

\clearpage
\subsection{Country-fixed-effects prediction models}

The regression exercise estimates linear probability models with country fixed
effects and Driscoll--Kraay covariance estimates. Across horizons, high debt
growth and the R-Zone interaction carry the economically important effects;
price growth alone is less stable. Coefficients are not expected to match
exactly because R-Zone assignments are sensitive to the unavailable
proprietary equity histories.

\begin{landscape}
\captionof{table}{\textbf{Replicated Table 4: crisis prediction by sector.}
Coefficients are percentage-point changes in the probability of a crisis
within the stated horizon. The sum rows give the predicted effect when all
included high-growth indicators equal one. The table supports the paper's
qualitative conclusion that joint debt and asset-price booms are more
informative than price appreciation alone.}
\label{tab:table4-historical}
\input{../_output/table4_replication.tex}
\end{landscape}

\begin{figure}[H]
\centering
\includegraphics[width=0.90\textwidth]{../_output/figure3_fraction_countries_rzone.jpg}
\caption{\textbf{Replicated Figure 3: fraction of countries in an R-Zone,
1950--2012.} The global share is the percentage of available countries in each
sector's R-Zone. The figure shows that warning episodes are synchronized across
countries in a few pronounced waves, especially around the late 1980s and
mid-2000s, rather than being evenly distributed through time.}
\label{fig:figure3-historical}
\end{figure}

\begin{table}[H]
\centering
\caption{\textbf{Numerical validation of the assigned historical exhibits.}
The validation suite compares every transcribed paper benchmark with the
recalculated value under a project-defined, paper-scaled tolerance. The bounds
use published dispersion, cell size, standard errors, counts, and sample size;
the table reports rather than conceals comparisons outside those bounds.}
\label{tab:validation}
\input{../_output/report_validation_summary.tex}
\end{table}

\section{Post-publication update}

The update is kept separate from the historical replication. Predictors and
R-Zone classifications extend through 2025. Because a horizon-$h$ forecast at
year $t$ needs crisis coverage through $t+h$, the latest valid forecast origins
are 2024, 2023, 2022, and 2021 at horizons one through four. This distinction
prevents late observations with incomplete future outcomes from being treated
as crisis-free.

\begin{table}[H]
\centering
\caption{\textbf{Updated Table 1: summary statistics through the latest
available year.} Predictor rows use complete sector pairs through 2025, while
legacy descriptive series retain their documented stopping dates. Compared
with the historical table, later observations reduce average debt growth and
change several expanded-sample quantiles. The R-Zone assignments nevertheless
continue to use the historical gates rather than these recalculated quantiles.}
\label{tab:table1-updated}
\small
\input{../_output/table1_post_publication.tex}
\end{table}

\begin{table}[H]
\centering
\caption{\textbf{Post-2016 R-Zone tracker by year.} Counts report the number of
countries with a usable sector pair and the number classified in each R-Zone.
The takeaway is that the warning indicator was uncommon in the recent panel
and absent in both sectors in 2024 and 2025, despite broad predictor coverage.}
\label{tab:tracker}
\small
\input{../_output/report_tracker_summary.tex}
\end{table}

\begin{figure}[p]
\centering
\includegraphics[width=0.96\textwidth]{../_output/post_publication_business_rzone_timeline.png}
\caption{\textbf{Updated Figure 1: event history through 2025.} Historical BVX
crises are joined to our extension of BVX's own criteria, and business
R-Zones use unchanged historical thresholds. The extension adds many predictor
observations but no confirmed post-2016 systemic onset inside the 42-country
sample.}
\label{fig:figure1-updated}
\end{figure}

\begin{figure}[p]
\ContinuedFloat
\centering
\includegraphics[width=0.96\textwidth]{../_output/post_publication_household_rzone_timeline.png}
\caption[]{\textbf{Updated Figure 1 continued: household event history through
2025.} Taken together, the updated panels are most informative about how
frequently the warning state has occurred rather than about newly realized
crisis performance.}
\end{figure}

\clearpage
\captionof{table}{\textbf{Updated Table 3: crisis probabilities using expanded
predictor and outcome coverage.} The table fully recalculates distributions,
frequencies, and differences from the median cell while retaining historical
quantile assignments. The updated upper-right cells remain riskier, but the
estimates partly dilute because the additional years contribute many
non-R-Zone observations and no newly classified systemic crisis in sample.}
\label{tab:table3-updated}
\input{../_output/table3_post_publication.tex}

\clearpage
\begin{landscape}
\captionof{table}{\textbf{Updated Table 4: re-estimated crisis-prediction
models.} Every coefficient is re-estimated rather than carried forward from
the historical table. Sample sizes decline with the horizon because longer
future windows require earlier forecast origins. The R-Zone effects remain
economically meaningful, although the update is not a strong out-of-sample
crisis test because no new systemic onset occurs within the country universe.}
\label{tab:table4-updated}
\input{../_output/table4_post_publication.tex}
\end{landscape}

\begin{figure}[H]
\centering
\includegraphics[width=0.90\textwidth]{../_output/figure3_fraction_countries_rzone_post_publication.jpg}
\caption{\textbf{Updated Figure 3: fraction of available countries in an
R-Zone through 2025.} The historical waves remain visible, while recent global
R-Zone shares are much lower. The important interpretation is descriptive:
the synchronized joint boom condition that characterized earlier peaks has not
been widespread in the most recent observations.}
\label{fig:figure3-updated}
\end{figure}

\clearpage
\section{Extensions and diagnostics}

\subsection{Bank-balance-sheet fragility}

The fragility extension adds high noncore funding, high loans-to-deposits, and
thin bank capital (the flipped lev inequality) one at a time, plus their
interaction with the sector R-Zone. All comparisons use the same complete-case sample as the corresponding
GHSS model. At the three-year horizon, noncore funding and loans-to-deposits
have positive standalone associations with crisis risk and improve within
fit. Interaction estimates are unstable; in particular, the signs vary across
measures and sectors, and the proposal's precommitted cutoff sweep shows why:
the joint R-Zone-and-fragile cell holds only a handful of country-years at
every candidate threshold, so no flag-based interaction estimate is well
identified.

\begin{table}[H]
\centering
\caption{\textbf{Three-year bank-fragility extension.} High fragility and
interaction coefficients are percentage points; $R^2_0$ and $R^2_1$ refer to
the same-sample GHSS and extended models. The table's main lesson is that
noncore funding and loans-to-deposits add standalone information, while the
interaction effect is too measure- and sector-sensitive to support a single
amplification mechanism.}
\label{tab:fragility}
\small
\input{../_output/report_fragility_summary.tex}

\medskip
\footnotesize
\textit{Notes:} *, **, and *** denote conventional significance at the 10\%,
5\%, and 1\% levels using Driscoll--Kraay p-values.
\end{table}

The proposal's continuous robustness check replaces each 0/1 flag with the
standardized ratio level (sign-flipped for lev, so a larger value always means
more fragile) and settles the question the flags cannot: with the full
variation in use, the standalone funding-fragility terms are positive and
significant while the R-Zone interactions remain indistinguishable from zero
for both funding measures. Funding fragility is an additive second crisis
channel, not an amplifier of credit booms, and adding it roughly doubles the
baseline model's within fit at the three- and four-year horizons.

\begin{table}[H]
\centering
\caption{\textbf{Three-year continuous fragility extension.} Coefficients are
percentage points of crisis probability per one standard deviation of the
ratio, standardized on the historical estimation sample. The table's main
lesson is that the funding-fragility level adds strong standalone signal and
roughly doubles within fit, while the R-Zone interaction stays a null once
the full variation replaces the thin flag cells: fragility adds to boom risk
rather than multiplying it.}
\label{tab:fragility-continuous}
\small
\input{../_output/report_fragility_continuous.tex}

\medskip
\footnotesize
\textit{Notes:} *, **, and *** denote conventional significance at the 10\%,
5\%, and 1\% levels using Driscoll--Kraay p-values.
\end{table}

Two further precommitted checks pin down what kind of predictor fragility is.
Racing the noncore level against its own three-year change shows the level
retains its full effect while the change carries nothing, so the channel is a
stock, the opposite time signature of the paper's boom variables, which are
all three-year flows. Adding a quadratic term leaves an economically and
statistically negligible coefficient at every horizon, so the dose-response
is linear: there is no threshold at which fragility switches from safe to
dangerous, which is why no flag version of this variable can achieve
R-Zone-like precision and why the monitor presents fragility as a continuous
gauge rather than a binary alarm.

\begin{table}[H]
\centering
\caption{\textbf{Functional form of the noncore fragility channel.} Level and
three-year-change coefficients come from their joint specification; the
quadratic term comes from the curvature check. Both checks were precommitted.
The takeaway is that fragility risk sits in the level, not the recent change,
and rises linearly with no threshold.}
\label{tab:fragility-form}
\small
\input{../_output/report_fragility_form_checks.tex}

\medskip
\footnotesize
\textit{Notes:} Coefficients are percentage points per standard deviation
(the quadratic per squared standard deviation). *, **, and *** denote
significance at the 10\%, 5\%, and 1\% levels using Driscoll--Kraay p-values.
\end{table}

\subsection{Crises missed by the R-Zone}

Of 54 historical crisis onsets in the diagnostic sample, 26 have no business
or household R-Zone in the preceding three years. Fragility data are available
for 16 of those missed crises. Seven have prior noncore funding at or above the
80th fragility percentile, six have high loans-to-deposits, and two have thin
bank capital.
Fragility therefore identifies a meaningful subset of misses but does not
explain all of them. Missing JST coverage is itself an important constraint.

\input{../_output/report_missed_crises.tex}

\subsection{Autoregressive comparison}

The dynamic specification adds current and lagged crisis states and lagged
predictors to the same-sample GHSS model. Within $R^2$ and in-sample AUC rise at
every horizon, most visibly at the one- and two-year horizons; crisis memory
sharpens the near view and adds little at the horizons the R-Zone is designed
for. Because the AUC values are in-sample, the exercise establishes
incremental fit rather than a definitive out-of-sample forecasting gain.

The paper's footnote 10 asserts, without a supporting table, that a local
projection with these dynamic terms gives ``qualitatively similar results.''
Our proposal pinned that phrase down before estimation: signs must match,
significance must survive at the same horizons, and the dynamic R-Zone
coefficient must sit within one standard error of its static value. All three
criteria hold in every sector-horizon cell, so the footnote's claim, which
the literature has so far taken on faith, verifies on our panel.

\begin{table}[H]
\centering
\caption{\textbf{The footnote-10 verdict.} Static and dynamic R-Zone
coefficients (percentage points) with the tightest precommitted criterion,
whether the dynamic estimate lies within one static standard error. Every
cell passes all three precommitted criteria, confirming the paper's unshown
robustness claim.}
\label{tab:fn10}
\small
\input{../_output/report_footnote10_verdict.tex}

\medskip
\footnotesize
\textit{Notes:} *, **, and *** denote significance at the 10\%, 5\%, and 1\%
levels using Driscoll--Kraay p-values.
\end{table}

\begin{table}[H]
\centering
\caption{\textbf{Static versus autoregressive crisis models.} The table
compares the R-Zone coefficient, within fit, and in-sample AUC on identical
samples. Adding memory consistently improves fit and classification, but the
gains are moderate at longer horizons and should not be interpreted as
out-of-sample evidence. Subscripts 0 and 1 denote the static and dynamic
models.}
\label{tab:dynamic}
\scriptsize
\input{../_output/report_dynamic_summary.tex}

\medskip
\footnotesize
\textit{Notes:} R-Zone coefficients are percentage points. *, **, and ***
denote significance at the 10\%, 5\%, and 1\% levels.
\end{table}

\subsection{Crisis risk after leaving the R-Zone}

The R-Zone's price condition switches off when a boom stalls, which is often
the prelude to the bust, so a country can leave the zone at the moment of
greatest danger. We test this with mutually exclusive one-, two-, and
three-years-since-exit indicators added to the same-sample baseline, with the
expectation of elevated post-exit risk stated before estimation. On all
origins the sectors split sharply: household exit-years carry elevated risk
at every horizon on top of the current-zone effect, while the business zone
shows nothing. That headline, however, does not survive its own robustness
check. A crisis itself crushes prices and ejects a country from the zone, so
roughly two-fifths of household exit-years sit within two years after a BVX
crisis, and the estimate leans on the 2008--2011 wave, where BVX record
second crisis onsets (Spain, Greece, Ireland, Denmark, Italy) that exit-years
from the first leg mechanically ``predict.'' Restricting to origins with no
crisis in the current or three prior years leaves the coefficients similar
but mostly removes their significance, and removing the five 2008-wave
countries' post-crisis origins eliminates the effect. We therefore report the
post-exit elevation as suggestive and concentrated in crisis waves, not as an
established warning property, and the honest design lesson for the monitor is
narrower: exit-years are not evidence of safety, but we cannot show they
carry standalone advance warning.

\begin{table}[H]
\centering
\caption{\textbf{Crisis risk after leaving the R-Zone.} Coefficients are
percentage points relative to country-years outside the zone with no recent
exit; the final column re-estimates the pooled exit effect on origins with no
BVX crisis in the current or three prior years. The takeaway is two-sided:
household post-exit risk looks elevated on all origins, but much of it
reflects overlap with ongoing crisis episodes, and the clean-origin estimates
are similar in size yet mostly insignificant.}
\label{tab:rzone-exit}
\small
\input{../_output/report_rzone_exit.tex}

\medskip
\footnotesize
\textit{Notes:} ``In zone'' and ``Recent exit'' come from the pooled
specification on all origins; the three exit-year columns come from the
mutually exclusive bucket specification; ``No-crisis exit'' is the pooled
exit coefficient on the no-recent-crisis sample. The estimation sample
excludes each country's earliest years, where exit status is undeterminable,
which raises the base crisis rate relative to Table~\ref{tab:table4-historical}'s
sample; the ``In zone'' column is therefore not comparable across tables.
*, **, and *** denote significance at the 10\%, 5\%, and 1\% levels using
Driscoll--Kraay p-values.
\end{table}

\subsection{United States and the 2023 banking stress}

The United States was not in either aggregate R-Zone from 2019 through 2023.
This is an informative failure mode rather than a reason to retroactively
redefine the model: the R-Zone is designed to detect borrower-side credit and
asset-price booms, whereas the March 2023 episode prominently involved
institution-specific duration, funding, uninsured-deposit, and run risk. The
public aggregate banking ratios are also unavailable after 2020 in this panel,
so the table cannot stand in for a bank-level SVB analysis.

\begin{table}[H]
\centering
\caption{\textbf{United States predictor case study, 2019--2023.} Three-year
business and household debt growth and real asset-price growth never jointly
cross their historical thresholds. The takeaway is that the original R-Zone
did not signal the 2023 stress, illustrating the difference between an
aggregate borrower boom and concentrated bank funding or duration risk.}
\label{tab:usa-case}
\small
\input{../_output/report_usa_case_study.tex}
\end{table}

\section{Assessment: successes, challenges, and lessons}

\subsection*{Where the replication succeeded}

The project delivers an end-to-end public-data reconstruction with a unique
country-year key, documented source priorities, reproducible transformations,
and automated validation. It closely recovers the paper's R-Zone frequencies,
preserves the joint debt--price crisis gradient, and produces regression
effects with the same broad economic interpretation. All 586 transcribed
historical benchmarks are audited under paper-scaled tolerances: 545 pass and
41 do not. Table 1 and Figure 3 pass completely, while the remaining misses are
concentrated in sparse Table 3 cells and Table 4 inference and fit statistics.
The automated suite verifies the main growth, threshold, outcome-window,
sample, and updated-coverage rules and marks the known exhibit-level misses as
expected failures. The post-publication work recalculates rather than merely
relabels the tables and figures.

\subsection*{Where the project faced challenges}

The largest obstacle is data equivalence. The paper partly relies on Global
Financial Data and Bloomberg equity histories, which are proprietary. IMF,
JST, and OECD substitutions are defensible and reproduce thresholds closely,
but different volatility and country coverage alter some individual R-Zone
assignments. Public data vintages have also been revised since the paper, and
long historical panels require careful source-by-source splicing. On inference,
the project implements Driscoll--Kraay covariance estimates but not the paper's
additional Kiefer--Vogelsang finite-sample correction.

The modern extension has a second limitation: the updated crisis chronology
contains no confirmed post-2016 onset within the 42-country sample. The added
years therefore improve knowledge about the prevalence of warning states but
provide little power for judging their realized post-publication crisis
performance. Likewise, JST aggregate banking variables cover far fewer
country-years and do not capture the institution-specific vulnerabilities that
matter in cases such as Continental Illinois or Silicon Valley Bank.

\subsection*{Overall conclusion}

The public reconstruction supports the paper's central qualitative result:
rapid debt growth is most concerning when accompanied by rapid real asset-price
appreciation. The R-Zone is a useful macro-financial warning state, not a
complete crisis model. Bank fragility and crisis memory add information on
complementary horizons, memory sharpening the near view and the
funding-fragility level, an additive and threshold-free channel, carrying the
three-to-four-year signal, while
missed crises and the 2023 United States episode show why aggregate
country-level boom indicators should be supplemented with banking-system and
institution-level measures. The updated data do not overturn the historical
relationship, but the absence of recent in-sample systemic onsets means that a
strong modern validation claim would be premature.

\appendix
\clearpage
\section{Detailed Table 1 divergence report}

\captionof{table}{\textbf{Historical Table 1 replication versus published
values.} This appendix makes the numerical differences transparent. The
largest substantive discrepancy is tied to the free equity-price replacement;
other moments and quantiles are generally close enough to support the main
replication exercise.}
\label{tab:table1-comparison}
\scriptsize
\input{../_output/table1_comparison.tex}

\section{References and source documentation}

Greenwood, Robin, Samuel G. Hanson, Andrei Shleifer, and Jakob Ahm
S{\o}rensen. 2022. ``Predictable Financial Crises.'' \textit{Journal of
Finance} 77(2).

Baron, Matthew, Emil Verner, and Wei Xiong. 2021. ``Banking Crises Without
Panics.'' \textit{Quarterly Journal of Economics} 136(1).

Additional public series are drawn from the IMF Global Debt Database and
former share-price statistics, the BIS Total Credit and Residential Property
Price databases, the JST Macrohistory database, OECD share-price and house-price
indicators, and World Bank World Development Indicators.

\end{document}
